%!TEX TS-program = Arara
% arara: pdflatex: { draft : yes }
% arara: nomencl
% arara: biber
% arara: pdflatex

%-------------------------------------------------------------------------------
%Dokumentenklasse und Schriftarten
%-------------------------------------------------------------------------------
\documentclass[12pt, a4paper, ngerman,toc=chapterentrydotfill, listof=totoc,bibliography=totoc,captions=figureheading,hidelinks,listof=nochaptergap,fleqn]{scrreprt} 
\usepackage[T1]{fontenc} 
\usepackage[utf8]{inputenc}

% \usepackage{arev} % Mathematik-Font ohne Serifen
% \usepackage{helvet} % Helvetica (Arial) ohne Serifen
% \renewcommand{\rmdefault}{\sfdefault} % Standardschrift setzen

\usepackage{microtype} % Mikrotypografie
\usepackage[ngerman]{babel}%Sprache und Silbentrennung
\usepackage{paralist}

\usepackage[onehalfspacing]{setspace}%Zeilenabstand auf 1,5 stellen
\usepackage[a4paper, left=4cm, right=2cm, top=1cm, bottom=2cm, headsep=2cm, includehead=true, footskip=0cm, includefoot=true, headheight = 1cm]{geometry} 


%-------------------------------------------------------------------------------
%Mathematische Ausdrücke
%-------------------------------------------------------------------------------
\usepackage{amsmath}
\usepackage{amssymb}
\counterwithout{equation}{chapter}%Nummerierung aufsteigend ohne Kapitelnummer

% Mathe Einrückung
\setlength{\mathindent}{1cm}

%-------------------------------------------------------------------------------
%Link Umbrueche im Verzeichnis
%-------------------------------------------------------------------------------
\usepackage{hyperref}%zum Setzen von PDF Daten und umbrechen von Links
\sloppy

%-------------------------------------------------------------------------------
%Meta Infos laden
%-------------------------------------------------------------------------------
%-----------------------------------
% Meta Informationen zur Arbeit
%-----------------------------------

% Autor
\newcommand{\myAutor}{Josh Arndt}
% Adresse
\newcommand{\myAdresse}{Zum Stadion 94 \\ \> \> 40764 Langenfeld}
% Titel der Arbeit
\newcommand{\myTitel}{LATEX-Vorlage - Fom Leitfaden}
% Betreuer
\newcommand{\myBetreuer}{Prof. Dr. Max Mustermann}
% Lehrveranstaltung
\newcommand{\myLehrveranstaltung}{Mustermodul}
% Matrikelnummer
\newcommand{\myMatrikelNr}{453673}
% Ort
\newcommand{\myOrt}{Köln}
% Datum der Abgabe
\newcommand{\myAbgabeDatum}{30. Januar 2019}
% Semesterzahl
\newcommand{\mySemesterZahl}{3}
% Name der Hochschule
\newcommand{\myHochschulName}{FOM Hochschule für Oekonomie \& Management}
% Standort der Hochschule
\newcommand{\myHochschulStandort}{Standort Köln}
% Studiengang
\newcommand{\myStudiengang}{Wirtschaftsinformatik}
% Art der Arbeit
\newcommand{\myThesisArt}{Seminararbeit}
% Zu erlangender akademische Grad
\newcommand{\myAkademischerGrad}{Bachelor of Science (B. Sc.)}
% Firma
\newcommand{\myFirma}{Mustermann GmbH}

%---------------------------------------
% PDF Meta setzen
%---------------------------------------
\usepackage{hyperref}
\hypersetup{
    pdfinfo={
        Title={\myTitel},
        Subject={\myStudiengang},
        Author={\myAutor},
        Build=1.1
    }
}

%-------------------------------------------------------------------------------
%Abbildungen einbinden
%-------------------------------------------------------------------------------
\usepackage{graphicx}
\usepackage{float}%Festlegen von Float Objects mit [H]
\usepackage{chngcntr}%Paket für Abbildungsnummerierung ändern
\counterwithout{footnote}{chapter}%ändern der Fussnoten-Nummerierung zu aufsteigend
\counterwithout{figure}{chapter}%ändern der Nummerierung zu aufsteigend
\counterwithout{table}{chapter}%ändern der Nummerierung zu aufsteigend
\usepackage[labelfont=bf,position=top, belowskip = 6pt, aboveskip=12pt]{caption}%Abstand zum Text festlegen

%-------------------------------------------------------------------------------
%Zitate und Literaturverzeichnis
%-------------------------------------------------------------------------------
\usepackage[babel,german=guillemets]{csquotes} %Festlegen der Anführungszeichen
\usepackage[backend=biber, style=ext-authoryear, isbn=false, maxcitenames=2, maxbibnames=999, mergedate=false, innamebeforetitle, dashed=false, autocite=footnote, mincrossrefs=1, autopunct=false, urldateusetime=true]{biblatex}%Bibliotheksverwaltung (siehe goLatex FOM)

\usepackage{blindtext}%Paket für Blindtext Generierung

%-------------------------------------------------------------------------------
%Abkuerzungsverzeichnis
%-------------------------------------------------------------------------------
\usepackage[intoc]{nomencl}
\renewcommand{\nomname}{Abkürzungsverzeichnis}%umbenennen
\setlength{\nomlabelwidth}{.15\textwidth}%Breite des Tabstops festlegen
%\renewcommand{\nomlabel}[1]{#1 \dotfill}%Füllen mit Punkt
\setlength{\nomitemsep}{-\parsep}
\makenomenclature

%-------------------------------------------------------------------------------
%Weitere Einstellungen
%-------------------------------------------------------------------------------

\setlength{\parindent}{0em}%Einrückung nach Absatz auf 0 setzen}%Seitenränder gesetzt
\setlength{\parskip}{6pt}%Absatzabstand
%-----------Kopf/Fußzeile--------------

\usepackage[automark, headsepline=.4pt]{scrlayer-scrpage}%Paket für Kopfzeile
\newpairofpagestyles{myheading}%Kopfzeile definieren
{
	\ihead{\leftmark}%Kapitelname links in Kopfzeile
	\chead{}%nichts in der Mitte
	\ohead{\thepage}%Seitenzahl rechts in Kopfzeile
}

\usepackage[automark, headsepline=.4pt]{scrlayer-scrpage}%Paket für Kopfzeile
\newpairofpagestyles{leer}%Kopfzeile definieren
{
	\ihead{}%Kapitelname links in Kopfzeile
	\chead{}%nichts in der Mitte
	\ohead{}%Seitenzahl rechts in Kopfzeile
}


\pagestyle{myheading}%Kopfzeile aktiv machen
\renewcommand*\chapterpagestyle{myheading}%Kopfzeile auch auf Seiten mit neuem Kapitel anzeigen

%----------Kapitelüberschriften formatieren-------

\setkomafont{disposition}{\rmfamily\bfseries}%Alle Überschriften mit Serifen und Fett
\setkomafont{chapter}{\huge}%Kapitelüberschriften in Groß
\setkomafont{section}{\large}%Sektionsüberschriften in Groß

\RedeclareSectionCommand[beforeskip=0pt, afterskip=6pt]{chapter}%12pt vor und 6pt nach Kapitelüberschrift
\RedeclareSectionCommand[beforeskip=12pt, afterskip=6pt]{section}%12pt vor und 6pt nach Kapitelüberschrift
%Datei mit Zitieranpassungen und Format des Literaturverzeichnisses
%https://golatex.de/eigener-zitierstil-anpassung-authoryear-icomp-t20497,start,30.html

\setlength\bibhang{1cm} 

\DefineBibliographyStrings{german}{% 
urlseen = {Zugriff am}, 
}

\renewcommand*{\finentrypunct}{} 

\renewcommand*{\newunitpunct}{\addcomma\space} 
\DeclareDelimFormat[bib,biblist]{nametitledelim}{\addcolon\space} 
\DeclareDelimFormat{titleyeardelim}{\newunitpunct} 
\renewcommand*{\mkbibnamefamily}{\mkbibemph} 
\renewcommand*{\mkbibnamegiven}{\mkbibemph} 
\renewcommand*{\mkbibnamesuffix}{\mkbibemph} 
\renewcommand*{\mkbibnameprefix}{\mkbibemph} 
\DeclareDelimAlias{finalnamedelim}{multinamedelim} 

\DeclareNameFormat{family-given:init}{% 
  \usebibmacro{name:family-given} 
    {\namepartfamily} 
    {\namepartgiveni} 
    {\namepartprefix} 
    {\namepartsuffix}% 
  \usebibmacro{name:andothers}} 

\DeclareNameAlias{default}{family-given} 
\DeclareNameAlias{sortname}{default} 
\DeclareNameAlias{labelname}{family-given:init}

\DeclareFieldFormat{editortype}{\mkbibparens{#1}} 
\DeclareDelimFormat{editortypedelim}{\addspace} 
\DeclareFieldFormat{translatortype}{\mkbibparens{#1}} 
\DeclareDelimFormat{translatortypedelim}{\addspace} 
\DeclareDelimFormat[bib,biblist]{innametitledelim}{\addcomma\space} 

\DeclareFieldFormat*{citetitle}{#1} 
\DeclareFieldFormat*{title}{#1} 
\DeclareFieldFormat*{booktitle}{#1} 
\DeclareFieldFormat*{journaltitle}{#1} 

\DeclareFieldFormat[online]{date}{\mkbibparens{#1}} 
\usepackage{xpatch} 
\xpatchbibdriver{online} 
  {\printlist{organization+location+date}} 
  {\setunit{\addspace}\usebibmacro{organization+location+date}} 
  {}{} 


\renewbibmacro*{date+extradate}{% 
  \printtext[parens]{% 
    \printfield{usera}% 
    \setunit{\printdelim{titleyeardelim}}% 
    \printlabeldateextra}} 

\DefineBibliographyStrings{german}{ 
  nodate    = {{}o.\adddot J\adddot}, 
  andothers = {et\addabbrvspace al\adddot} 
} 

\DeclareSourcemap{ 
  \maps[datatype=bibtex]{ 
    \map{ 
      \step[notfield=translator, final] 
      \step[notfield=editor, final] 
      \step[fieldset=author, fieldvalue={{{o\noexpand\adddot V\noexpand\adddot}}}] 
    } 
    \map{ 
      \pernottype{online} 
      \step[fieldset=location, fieldvalue={o\noexpand\adddot O\noexpand\adddot}] 
    } 
  } 
} 

\renewbibmacro*{cite}{% 
  \iffieldundef{shorthand} 
    {\ifthenelse{\ifnameundef{labelname}\OR\iffieldundef{labelyear}} 
       {\usebibmacro{cite:label}% 
        \setunit{\printdelim{nonametitledelim}}} 
       {\printnames{labelname}% 
        \setunit{\printdelim{nametitledelim}}}% 
     \printfield{usera}% 
     \setunit{\printdelim{titleyeardelim}}% 
     \usebibmacro{cite:labeldate+extradate}} 
    {\usebibmacro{cite:shorthand}}} 

\renewbibmacro*{postnote}{% 
  \setunit{\postnotedelim}% 
  \iffieldundef{postnote} 
    {\printtext{o.S\adddot}} 
    {\printfield{postnote}}} %Datei mit Zitieranpassungen und Format des Literaturverzeichnisses
% sauber formatierter Quelltext
\usepackage{listings}
\lstset{numbers=left,
	numberstyle=\tiny,
	numbersep=5pt,
	breaklines=true,
	showstringspaces=false,
	frame=l ,
	xleftmargin=5pt,
	xrightmargin=5pt,
	basicstyle=\ttfamily\scriptsize,
	stepnumber=1,
	keywordstyle=\color{blue},          % keyword style
  	commentstyle=\color{dkgreen},       % comment style
  	stringstyle=\color{mauve}         % string literal style
}%Sauber formatierter Quelltext
\addbibresource{lit.bib}%include der Literatur Datenbank
%\nocite{*}%UNBEDINGT AUSKOMMENTIEREN. DURCH DIESEN BEFEHL WERDEN AUCH NICHT ZITIERTE QUELLEN VERZEICHNET
%\addbibresource{biblatex-examples.bib}%Beispielliteratur einfügen



%-------------------------------------------------------------------------------
%Dokumentbeginn
%-------------------------------------------------------------------------------
\begin{document}%Beginn des Dokumentes
\newcounter{savepage}
\pagenumbering{Roman}%Die ersten Seiten römisch nummerieren

%-------------------------------------------------------------------------------
%Titelblatt
%-------------------------------------------------------------------------------
\begin{titlepage}
	\newgeometry{left=4cm, right=2cm, top=4cm, bottom=2cm}
	
		\noindent \textbf{\myHochschulName~Essen}\\
		\noindent \textbf{\myHochschulStandort}\\
		\vspace{1cm}\\
		\noindent Berufsbegleitender Studiengang zum\\
		\noindent\myAkademischerGrad~\myStudiengang \\ 
		\vspace{1cm} \\		
		\mySemesterZahl. Semester\\
		\vspace{0.5cm}\\
		%\textbf{zur Erlangung des Grades eines}\\
		%\textbf{\myAkademischerGrad}\\
		\textbf{Seminararbeit in \myLehrveranstaltung}\\
		\vspace{2cm}\\
		\Huge{\myTitel}\\
		\vspace{0.2cm}
	\normalsize
	\vfill
	\begin{tabbing}
		Linkslinks \= Mitte \= Rechts\kill
		Betreuer: \> \> \myBetreuer\\
		Autor: \> \> \myAutor\\
		Matrikelnr.: \> \>   \myMatrikelNr\\
		Abgabedatum: \> \> \myAbgabeDatum
	\end{tabbing}
\end{titlepage}

%-------------------------------------------------------------------------------
%Verzeichnisse
%-------------------------------------------------------------------------------
\setcounter{page}{2}%TOC als erste Seite
\tableofcontents%Inhaltsverzeichnis einfügen
\listoffigures%Abbildungsverzeichnis einfügen
\listoftables%Tabellenverzeichnis einfügen
\printnomenclature%Abkürzungsverzeichnis einfügen
\cleardoublepage

%-------------------------------------------------------------------------------
%Text
%-------------------------------------------------------------------------------
\setcounter{savepage}{\arabic{page}}%Seitenstand der römischen Nummerierung speichern um später dort weiterzuzählen
\pagenumbering{arabic}%arabische Nummerierung

\chapter{Einleitung}
\label{ch:einleitung}

\section{Umfeld}
Diese \LaTeX -Vorlage ist zur Nutzung bei der Erstellung von Hausarbeiten und der Bachelor-/Master-Thesis vorgesehen. Dabei wurde sich auf den Leitfaden zur Erstellung wissenschaftlicher Arbeiten aus dem Februar 2018 bezogen.\autocite{leitfaden}

\section{Motivation}
Gerade in den letzten Jahren wurde der Leitfaden zur Erstellung wissenschaftlicher Arbeiten an der FOM (\nomenclature{FOM}{Hochschule für Oekonomie und Management} Hochschule für Oekonomie und Management) immer komplexer. Viele Vorschriften lassen sich nicht wie früher noch mit den \LaTeX\-Standardeinstellungen oder kleineren Anpassungen erreichen, sondern müssen, wie beispielsweise das Zitieren (siehe Kapitel \ref{ch:zitieren}, S. \pageref{ch:zitieren}).

Um das Anfertigen von Hausarbeiten etc. zu vereinfachen, wurde dieses Template angefangen. Es soll den Kommilitonen als Anhaltspunkt für neue Arbeiten dienen und durch gemeinsame Weiterarbeit zunehmend wachsen. Aufgrunddessen wird selbstverständlich in keinster Weise die Richtigkeit der Formatierung garantiert. Ich habe mir zwar größte Mühe gegeben, aber Fehler passieren immer. Bei gefundenen Fehlern würde ich mich sehr über Rückmeldung freuen. Der Fehler wird dann schnellstmöglich behoben und schon haben wir alle etwas davon. :-)

Das folgende Dokument erklärt einige Funktionen des Templates, und wie ihr diese mit möglichst wenig Aufwand nutzen könnt. Auch hier stehe ich (oder ggf. andere Nutzer des Templates) bestimmt für Fragen zur Verfügung.

\section{Abgrenzung}
In der Anleitung werde ich lediglich auf die Nutzung auf Windows-Geräten eingehen. Zwar wird sich die Syntax innerhalb der .tex Dateien nicht ändern, doch kann es sein, dass die erstmalige Installation sich unterscheidet.
\chapter{Kapitel1}
\blindtext{}\autocite{elk}

\begin{figure}[!htb]
  \centering
  \caption{FOM Logo 2019}
  \includegraphics[width=8cm]{Abbildungen/fom_logo.png} 
  \\Quelle: \cite[12]{becker2016data}
  \label{fom_logo}
\end{figure}

\blindtext{}
\chapter{Schluss}
Hier würde der Schluss bzw. das Fazit der Arbeit eingefügt werden.%Datei mit eigentlichem Text einbinden

% \chapter{Anhang}
\renewcommand\thesection{Anhang \arabic{section}:} %https://tex.stackexchange.com/questions/3177/how-to-change-the-numbering-of-part-chapter-section-to-alphabetical-r
\addtocontents{toc}{\protect\setcounter{tocdepth}{-1}}
%-------------------------------------------------------------------------------

\begin{minipage}{\textwidth}
	\section{Versuch Nummer Eins}
	\label{anh:versuch_nummer_eins}
	\includegraphics[width=0.5\linewidth]{Abbildungen/fom_logo.png} 
	\par\smallskip
	Quelle: \cite{logo}
\end{minipage}

\addtocontents{toc}{\protect\setcounter{tocdepth}{1}}
\renewcommand\thesection{\thechapter.\arabic{section}} Anhänge
%-------------------------------------------------------------------------------
%Verzeichnisse
%-------------------------------------------------------------------------------
\pagenumbering{Roman}%römische Nummerierung
\setcounter{page}{\thesavepage}%Seitencounter auf die gespeicherte Zahl setzen

%--------FÜR HAUSARBEIT NACH LEITFADEN----------------------
\addchap{Literaturverzeichnis}
\printbibliography[nottype=online,heading=none]%Literaturverzeichnis der NICHT-Onlinequellen drucken und ins Inhaltsverzeichnis einfügen
\section*{Internetquellen}
\printbibliography[type=online,heading=none,title={Internetquellen}]%Literaturverzeichnis der Onlinequellen drucken.
%--------FÜR HAUSARBEIT BEI KRAUSE----------------------
% \addchap{Literaturverzeichnis}
% \printbibliography[nottype=online, nottype=article,heading=none]%Literaturverzeichnis der NICHT-Artikel/Onlinequellen drucken.
% \addchap{Quellenverzeichnis}
% \printbibliography[filter=articleonline,heading=none]%Literaturverzeichnis der /ArtikelOnlinequellen drucken.
%-------------------------------------------------------

% % !TeX root = ../../main.tex

%-----------------------------------
% Ehrenwörtliche Erklärung
%-----------------------------------
\chapter*{Ehrenwörtliche Erklärung}
\pagenumbering{gobble} % Keine Seitenzahlen mehr
\renewcommand*\chapterpagestyle{leer}
Hiermit versichere ich, dass die vorliegende Arbeit von mir selbstständig und ohne unerlaubte Hilfe angefertigt worden ist, insbesondere dass ich alle Stellen, die wörtlich oder annähernd wörtlich aus Veröffentlichungen entnommen sind, durch Zitate als solche gekennzeichnet habe. Ich versichere auch, dass die von mir eingereichte schriftliche Version mit der digitalen Version übereinstimmt. Weiterhin erkläre ich, dass die Arbeit in gleicher oder ähnlicher Form noch keiner Prüfungsbehörde/Prüfungsstelle vorgelegen hat. Ich erkläre mich damit einverstanden, dass die Arbeit der Öffentlichkeit zugänglich gemacht wird. Ich erkläre mich damit einverstanden, dass die Digitalversion dieser Arbeit zwecks Plagiatsprüfung auf die Server externer Anbieter hoch geladen werden darf. Die Plagiatsprüfung stellt keine Zurverfügungstellung für die Öffentlichkeit dar.

\par\medskip
\par\medskip

\vspace{5cm}

\begin{table}[H]
	\centering
	\begin{tabular*}{\textwidth}{c @{\extracolsep{\fill}} ccccc}
		\myOrt, \today
		&
		% Hinterlege deine eingescannte Unterschrift im Verzeichnis /abbildungen und nenne sie unterschrift.png
		% Bilder mit transparentem Hintergrund können teils zu Problemen führen
		\includegraphics[width=0.35\textwidth]{Abbildungen/unterschrift.png}\vspace*{-0.35cm}
		\\
		\rule[0.5ex]{12em}{0.55pt} & \rule[0.5ex]{12em}{0.55pt} \\
		(Ort, Datum) & (Eigenhändige Unterschrift)
		\\
	\end{tabular*} \\
\end{table} Erklärung
\end{document}%Ende des Dokumentes
